% !TEX root = ../main.tex
\clearpage
\phantomsection
\addcontentsline{toc}{chapter}{Referencias}
\begin{thebibliography}{99}

\bibitem{iso29148}
ISO/IEC/IEEE,
\emph{Systems and software engineering---Life cycle processes---Requirements engineering},
ISO/IEC/IEEE 29148:2018, 2nd ed., 2018.
doi: 10.1109/IEEESTD.2018.8559686.

\bibitem{iso42010}
ISO/IEC/IEEE,
\emph{Software, systems and enterprise---Architecture description},
ISO/IEC/IEEE 42010:2022, 2022.

\bibitem{iso15939}
ISO/IEC/IEEE,
\emph{Systems and software engineering---Measurement process},
ISO/IEC/IEEE 15939:2017, 2017.

\bibitem{iso25010}
ISO/IEC,
\emph{Systems and software engineering---Systems and software Quality Requirements and Evaluation (SQuaRE)---Product quality model},
ISO/IEC 25010:2023, 2023.

\bibitem{iso25040}
ISO/IEC,
\emph{Systems and software engineering---Systems and software Quality Requirements and Evaluation (SQuaRE)---Quality evaluation framework},
ISO/IEC 25040:2024, 2024.

\bibitem{spdp2025}
Superintendencia de Protección de Datos Personales del Ecuador,
\emph{Resolución SPDP-SPD-2025-0006-R: cláusulas de protección de datos personales para contratos que incluyan transferencia o comunicación de datos personales},
6 de mayo de 2025.
Disponible en: \url{https://spdp.gob.ec/wp-content/uploads/2025/05/06.01.01-SPDP-SPD-2025-0006-R-clausulas-para-contratos-en-Ecuador-signed.pdf}.

\bibitem{sommerville2011}
I. Sommerville,
\emph{Software Engineering}, 9th ed.
Boston, MA: Addison-Wesley, 2011.

\bibitem{mavin2009}
A. Mavin, P. Wilkinson, A. Harwood, and M. Novak,
“Easy Approach to Requirements Syntax (EARS),”
in \emph{Proceedings of the 17th IEEE International Requirements Engineering Conference},
2009, pp. 317--322.
doi: 10.1109/RE.2009.9.

\bibitem{wagner2019}
S. Wagner et al.,
“Status Quo in Requirements Engineering: A Theory and a Global Family of Surveys,”
\emph{ACM Transactions on Software Engineering and Methodology},
vol. 28, no. 2, art. 9, 2019.
doi: 10.1145/3306607.

\bibitem{palomares2021}
C. Palomares, X. Franch, C. Quer, P. Chatzipetrou, L. López, and T. Gorschek,
“The state-of-practice in requirements elicitation: an extended interview study at 12 companies,”
\emph{Requirements Engineering}, vol. 26, pp. 273--299, 2021.
doi: 10.1007/s00766-020-00345-x.

\bibitem{mucha2024}
J. Mucha, A. Kaufmann, and D. Riehle,
“A systematic literature review of pre-requirements specification traceability,”
\emph{Requirements Engineering}, vol. 29, pp. 119--141, 2024.
doi: 10.1007/s00766-023-00412-z.

\bibitem{guo2025}
J. L. C. Guo, J.-P. Steghöfer, A. Vogelsang, and J. Cleland-Huang,
“Natural Language Processing for Requirements Traceability,”
in \emph{Handbook on Natural Language Processing for Requirements Engineering}.
Springer, 2025, pp. 89--116.
doi: 10.1007/978-3-031-73143-3\_4.

\bibitem{lewis2020}
P. Lewis et al.,
“Retrieval-Augmented Generation for Knowledge-Intensive NLP Tasks,”
in \emph{Advances in Neural Information Processing Systems 33},
2020, pp. 9459--9474.

\bibitem{reimers2019}
N. Reimers and I. Gurevych,
“Sentence-BERT: Sentence Embeddings using Siamese BERT-Networks,”
in \emph{Proceedings of the 2019 Conference on Empirical Methods in Natural Language Processing and the 9th International Joint Conference on Natural Language Processing},
2019, pp. 3982--3992.
doi: 10.18653/v1/D19-1410.

\bibitem{cormack2009}
G. V. Cormack, C. L. A. Clarke, and S. Buettcher,
“Reciprocal rank fusion outperforms Condorcet and individual rank learning methods,”
in \emph{Proceedings of the 32nd International ACM SIGIR Conference on Research and Development in Information Retrieval},
2009, pp. 758--759.
doi: 10.1145/1571941.1572114.

\bibitem{liu2024}
N. F. Liu et al.,
“Lost in the Middle: How Language Models Use Long Contexts,”
\emph{Transactions of the Association for Computational Linguistics},
vol. 12, pp. 157--173, 2024.
doi: 10.1162/tacl\_a\_00638.

\bibitem{norheim2024}
J. J. Norheim, E. Rebentisch, D. Xiao, L. Draeger, A. Kerbrat, and O. L. de Weck,
“Challenges in applying large language models to requirements engineering tasks,”
\emph{Design Science}, vol. 10, e16, 2024.
doi: 10.1017/dsj.2024.8.

\bibitem{hemmat2025}
A. Hemmat, M. Sharbaf, S. Kolahdouz-Rahimi, K. Lano, and S. Y. Tehrani,
“Research directions for using LLM in software requirement engineering: a systematic review,”
\emph{Frontiers in Computer Science}, vol. 7, 2025.
doi: 10.3389/fcomp.2025.1519437.

\bibitem{cheng2026}
H. Cheng, J. H. Husen, Y. Lu, T. Racharak, N. Yoshioka, N. Ubayashi, and H. Washizaki,
“Generative AI for Requirements Engineering: A Systematic Literature Review,”
\emph{Software: Practice and Experience}, vol. 56, no. 2, pp. 141--170, 2026.
doi: 10.1002/spe.70029.

\bibitem{sabetzadeh2025}
M. Sabetzadeh and C. Arora,
“Practical Guidelines for the Selection and Evaluation of Natural Language Processing Techniques in Requirements Engineering,”
in \emph{Handbook on Natural Language Processing for Requirements Engineering}.
Springer, 2025.

\bibitem{vogelsang2025}
A. Vogelsang and J. Fischbach,
“Using Large Language Models for Natural Language Processing Tasks in Requirements Engineering: A Systematic Guideline,”
in \emph{Handbook on Natural Language Processing for Requirements Engineering}.
Springer, 2025, pp. 435--456.
doi: 10.1007/978-3-031-73143-3\_16.

\bibitem{mare2024}
D. Jin, Z. Jin, X. Chen, and C. Wang,
“MARE: Multi-Agents Collaboration Framework for Requirements Engineering,”
arXiv:2405.03256, 2024.

\bibitem{sami2024}
M. A. Sami, M. Waseem, Z. Zhang, Z. Rasheed, K. Systä, and P. Abrahamsson,
“Early Results of an AI Multiagent System for Requirements Elicitation and Analysis,”
in \emph{Product-Focused Software Process Improvement},
2024, pp. 307--316.
doi: 10.1007/978-3-031-78386-9\_20.

\bibitem{iredev2025}
D. Jin, W. Sun, J. Huang, P. Liang, J. Xuan, Y. Liu, and Z. Jin,
“iReDev: A Knowledge-Driven Multi-Agent Framework for Intelligent Requirements Development,”
arXiv:2507.13081, 2025.

\bibitem{oriol2025}
M. Oriol, Q. Motger, J. Marco, and X. Franch,
“Multi-Agent Debate Strategies to Enhance Requirements Engineering with Large Language Models,”
in \emph{Proceedings of the IEEE 33rd International Requirements Engineering Conference},
2025, pp. 527--534.
doi: 10.1109/RE63999.2025.00063.

\bibitem{reqnet2026}
S. Saleem, M. N. Asim, and A. Dengel,
“ReqNet: an LLM-driven computational framework for automated requirements extraction from unstructured documents,”
\emph{Complex \& Intelligent Systems}, vol. 12, art. 38, 2026.
doi: 10.1007/s40747-025-02143-w.

\bibitem{paiva2026}
G. Paiva, E. D. Canedo, and G. P. Rocha Filho,
“From issue titles to requirements: an empirical study of large language models and prompt engineering strategies,”
\emph{Requirements Engineering}, vol. 31, art. 7, 2026.
doi: 10.1007/s00766-026-00462-z.

\bibitem{quare2026}
H. Cheng et al.,
“QUARE: Multi-Agent Negotiation for Balancing Quality Attributes in Requirements Engineering,”
arXiv:2603.11890, 2026.

\bibitem{tracellm2026}
N. Alturayeif, I. Ahmad, and J. Hassine,
“TraceLLM: leveraging large language models with prompt engineering for enhanced requirements traceability,”
\emph{Requirements Engineering}, vol. 31, no. 1, art. 6, 2026.
doi: 10.1007/s00766-026-00460-1.

\bibitem{ali2024}
S. J. Ali, V. Naganathan, and D. Bork,
“Establishing Traceability between Natural Language Requirements and Software Artifacts by Combining RAG and LLMs,”
in \emph{Conceptual Modeling},
2024, pp. 295--314.
doi: 10.1007/978-3-031-75872-0\_16.

\bibitem{lucassen2016}
G. Lucassen, F. Dalpiaz, J. M. E. M. van der Werf, and S. Brinkkemper,
“Improving agile requirements: the Quality User Story framework and tool,”
\emph{Requirements Engineering}, vol. 21, pp. 383--403, 2016.
doi: 10.1007/s00766-016-0250-x.

\bibitem{metagpt2024}
S. Hong et al.,
“MetaGPT: Meta Programming for a Multi-Agent Collaborative Framework,”
in \emph{International Conference on Learning Representations}, 2024.

\bibitem{wohlin2012}
C. Wohlin et al.,
\emph{Experimentation in Software Engineering}.
Berlin: Springer, 2012.
doi: 10.1007/978-3-642-29044-2.

\bibitem{berry2024}
D. M. Berry,
“Evaluation of Tools for Hairy Requirements Engineering and Software Engineering Tasks,”
in \emph{Proceedings of the IEEE International Requirements Engineering Conference},
2024.

\end{thebibliography}
